% Part: lambda-calculus
% Chapter: lambda-definability
% Section: introduction

\documentclass[../../../include/open-logic-section]{subfiles}

\begin{document}

\olfileid{lam}{ldf}{int}
\olsection{جان پہچان}

پہلی نظر وچ لامبڈا کیلکولس بس عبارتاں دا بڑا تجریدی حسابی
نظام لگدا اے۔ ایہ عبارتاں فنکشوناں تے اوہناں نوں ہور فنکشوناں
اُتے لاؤن دا اظہار کردیاں نیں۔ لامبڈا کیلکولس دی نحو وچ کوئی
ایہو جیہی گل نہیں جیہڑی دسے کہ ایہ کسے خاص قسم دیاں شَواں
اُتے فنکشن نیں۔ خاص کر کے، نحو وچ قدرتی عدداں دا کوئی ذکر
نہیں۔ ایہدیاں بُنیادی کارروائیاں، اطلاق تے لامبڈا تجرید، ہر
فنکشن لئی ورتیاں جا سکدیاں نیں، صرف قدرتی عدداں اُتے فنکشوناں
لئی نہیں۔

فیر وی تھوڑی جیہی سُوجھ نال لامبڈا کیلکولس وچ حسابی فنکشن،
یعنی قدرتی عدداں اُتے فنکشن، متعین کیتے جا سکدے نیں۔ ایس
لئی اسیں ہر قدرتی عدد~$n \in
\Nat$ لئی اک خاص $\lambd$-ٹرم~$\num n$ متعین کردے آں، جیہڑا
~$n$ دا \emph{چرچ عدد} اے۔ (چرچ عدداں دا ناں الونزو چرچ توں آیا اے۔)

\begin{defn}
  جے $n \in \Nat$، تاں اوہدے نال جُڑیا \emph{چرچ عدد} $\num{n}$
  عدد~$n$ نوں ظاہر کردا اے۔
  \[
    \num{n} \ident \lambd[fx][f^n(x)]
  \]
  ایتھے $f^n(x)$ اوہ نتیجہ اے جیہڑا $f$ نوں~$x$ اُتے $n$
  واری لاؤن نال ملدا اے۔ مثلاً $\num{0}$، $\lambd[fx][x]$ اے،
  تے $\num{3}$، $\lambd[fx][f(f(f\,x))]$ اے۔
\end{defn}

چرچ عدد $\num n$ نوں لامبڈا ٹرم دے روپ وچ لکھیا جاندا اے۔
اوہ ٹرم اک فنکشن ظاہر کردا اے جیہڑا دو دلیلاں $f$ تے $x$
قبول کردا اے تے $f^n(x)$ واپس دیندا اے۔ صاف اے کہ چرچ
عدد نارمل روپ وچ ہوندے نیں۔

قدرتی عدداں نوں لامبڈا کیلکولس وچ ظاہر کرنا اُدوں ای کم آؤندا
اے جدوں اسیں اوہناں نال حساب کر سکیے۔ لامبڈا کیلکولس وچ
چرچ عدداں نال حساب کرن دا مطلب اے کہ اک $\lambd$-ٹرم~$F$
ایہو جیہے چرچ عدد اُتے لائیے تے سانجھے ٹرم~$F\, \num n$ دی
تقلیل نارمل روپ تک کریے۔ جے ایہ ہمیشہ نارمل روپ تک پہنچدا
اے تے اوہ نارمل روپ ہمیشہ کوئی چرچ عدد~$\num m$ ہوندا اے،
تاں اسیں عدد~$m$ نوں حساب دی حاصل قدر سمجھ سکدے آں۔
فیر اسیں سمجھ سکدے آں کہ~$F$ اوہ فنکشن $f\colon \Nat \to \Nat$
متعین کردا اے جیہدے لئی $f(n) = m$ اُدوں تے صرف اُدوں
اے جدوں $F\, \num n \red \num m$۔ چرچ–روسر خاصیت کر کے
جے نارمل روپ موجود ہووے، تاں اکّو ہوندا اے۔ سو جے
$F\, \num n \red \num m$، تاں کوئی ہور نارمل ٹرم، خاص کر کے کوئی
ہور چرچ عدد، نہیں ہو سکدا جیہدے تک $F \, \num n$ دی تقلیل ہووے۔

اُلٹے پاسے، فنکشن $f\colon \Nat \to \Nat$ دِتّا ہووے، تاں
اسیں پُچھ سکدے آں کہ کوئی ٹرم $F$ موجود اے جیہڑا~$f$ نوں
ایس طرح متعین کرے۔ اوہ صورت وچ اسیں آکھدے آں کہ $F$،
~$f$ نوں \emph{!!{lambda define}s} تے $f$، !!{lambda definable}
اے۔ ایہ گل کئی دلیلاں والے تے جزوی فنکشوناں لئی وی عام کر سکدے آں۔

\begin{defn}
فرض کرو $f\colon \Nat^k \to \Nat$۔ اسیں آکھدے آں کہ
لامبڈا ٹرم~$F$، $f$ نوں \emph{!!{lambda define}s} جے ہر
$n_0$، \dots،~$n_{k-1}$ لئی
\[
F \, \num{n_0} \, \num{n_1} \dots \num{n_{k-1}} \red \num{f(n_0, n_1, \dots,
  n_{k-1})}
\]
ہووے جدوں $f(n_0, \dots, n_{k-1})$ معین ہووے، تے ہور صورت
وچ $F \, \num{n_0} \,
\num{n_1} \dots \num{n_{k-1}}$ دا کوئی نارمل روپ نہ ہووے۔
\end{defn}

ثابت قدر والے فنکشن بڑی سادی مثال نیں۔ ٹرم $C_k \ident
\lambd[x][\num{k}]$ اوہ فنکشن $c_k\colon \Nat \to
\Nat$ نوں !!{lambda define}s جیہدے لئی $c(n) = k$۔ کیوں جے $C_k \, \num n \ident
(\lambd[x][\num{k}])\num n \redone \num{k}$، بھاویں~$n$ کوئی
وی ہووے۔ عین اوہی قدر واپس دین والا فنکشن !!{lambda defined}
اے $\lambd[x][x]$ نال۔ ہور گُنجھلدار فنکشوناں نوں متعین کرنا
اوکھا ہوندا اے تے ایس لئی اکثر بڑی سُوجھ دی لوڑ پیندی اے۔
سو ایہ گل شاید حیرانی والی اے کہ ہر قابلِ حساب فنکشن
!!{lambda definable} اے۔ اُلٹی گل وی سچ اے۔ جے کوئی فنکشن
!!{lambda definable} اے، تاں اوہ قابلِ حساب اے۔

\end{document}
